\scanpage{15}
auquel on se ramène) on ne peut plus affirmer que $\Lambda$ et $\Lambda'$ sont isomorphes \deleted{\unsure{ce}} (ni $=$ loc. isom.); c'est \inserted{tjrs vrai (pr un \deleted{$\unsure{\mathscr P|P},\mathscr P'\ (M,N,\sigma)\to(M',N',\sigma')$ fixés)} $\Lambda\colon\mathscr P\to\mathscr P'$ fixé)} \deleted{\unsure{sauf}} \unsure{sss|ssi|si} $\operatorname{Ext}^1(S;M,N')=0$ (resp. $\underline{\operatorname{Ext}}^1(M,N')=0$). \deleted{Mais si on se donne}\inserted{On ne peut} affirmer non plus que tt hom $\lambda$ de $(M,N,\sigma)$ ds $(M',N',\sigma')$ provienne d'un $\Lambda$ \inserted{($=$ loc.)} \inserted{[c'est OK ssi $\operatorname{Ext}^{\unsure{2|1}}(S;\unsure{M},N')=0$ resp. $\underline{\operatorname{Ext}}^2(M,N')=0$]} (ni que tt $(M,N,\sigma)$ provienne d'un $\mathscr P$ -- cf. plus bas --)
\margin{On peut prendre pour $\lambda$ l'identité}
mais pour $\mathscr P,\mathscr P',\lambda$ fixé, le\deleted{\unsure{s}} \deleted{cat. des} champ des $\otimes$-hom AUC $\mathscr P\to\mathscr P'$ correspondant est un 1-\inserted{ps.}torseur sous le champ de Picard \deleted{des}\inserted{strict} d'invariants $\underline{\operatorname{Ext}}^1(M,N'),\underline{\operatorname{Hom}}(M,N')$ obtenu par tronçage de $R\underline{\operatorname{Hom}}(M,N')$; l'obstruction à trouver un tel \deleted{\unsure{hom}}\inserted{$\Lambda$} est dans $\operatorname{Ext}^2(S,M,N')$ [et à le trouver loc est dans $H^0(S,\underline{\operatorname{Ext}}^2(M,N'))$].

Tout ceci provient du jeu habituel à la Baer, et du fait (qui en résulte) que pour $M,N,\sigma$ \deleted{\unsure{$(M',N',\sigma')$}} fixés, les champs de Picard d'invariants
% archive sheet number: 14
